%! 3 by 3 Determinants — Unit 5, Lesson 5.1 — Lesson Plan
\documentclass[10pt]{article}
\usepackage{linalg-article}
\usepackage{linalg-boxes}
\usepackage{graphicx}   % -article does NOT load graphicx; needed for thumbnails/screenshots
\graphicspath{{images/}}
\newcommand{\TallMath}[1]{$\displaystyle #1\rule[-1.4em]{0pt}{3.2em}$}
\newcommand{\avec}{\mathbf{a}}
\newcommand{\bb}{\mathbf{b}}
\newcommand{\UnitNumberName}{Unit 5: Determinants and Linear Transformations \quad}
\newcommand{\LessonNumberName}{Lesson 5.1: 3 by 3 Determinants}

\begin{document}
\begin{center}
    {\Large\bfseries \CourseName: \SchoolYear \\
    {\small\bfseries \UnitNumberName \LessonNumberName}}
\end{center}

\begin{tcolorbox}[colback=blush, colframe=burgundy, boxrule=0.9pt, arc=2mm,
    left=2mm, right=2mm, top=1.5mm, bottom=1.5mm]
    \textbf{Primary Objective:} Students can compute a $3\times3$ determinant by \emph{cofactor
    expansion} --- expanding along a row into three $2\times2$ determinants (which they own from
    Lesson~5.0) with the sign pattern $+\,-\,+$ --- and read it \emph{geometrically}: $|\det A|$ is the
    \emph{volume} of the box (parallelepiped) the three columns span, and the \emph{sign} reports
    orientation. They can use the \emph{triangular shortcut} (determinant $=$ product of the diagonal)
    and explain why $\det=0$ means the three columns lie in \emph{one plane} --- the box collapses flat,
    and the matrix is \emph{not invertible} (the Unit~2 callback). They can say where Unit~5 goes next:
    properties of determinants (5.2) and transformations, where $|\det A|$ is the \emph{volume-scaling
    factor} (5.3).
\end{tcolorbox}

\renewcommand{\arraystretch}{1.4}
\begin{skillbox}[Priority Ideas \& Skills]{goldbox}
\begin{tabularx}{\linewidth}{@{}X|X@{}}
\textbf{Priority Ideas \& Skills} & \textbf{Key Understandings (the ``why'')} \\[2pt]
\hline
\vspace{-6pt}
\begin{itemize}[leftmargin=1.1em, nosep, after=\vspace{2pt}]
  \item Compute a $3\times3$ determinant by \textbf{cofactor expansion} ($+\,-\,+$).
  \item Use \textbf{minors}: cross out a row \& column to get each $2\times2$.
  \item Read $|\det A|$ as the \textbf{volume} of the box the columns span; use the \textbf{triangular} shortcut.
  \item Use $\det=0$: columns coplanar $\Leftrightarrow$ flat box $\Leftrightarrow$ \emph{not invertible}.
\end{itemize}
&
\vspace{-6pt}
\begin{itemize}[leftmargin=1.1em, nosep, after=\vspace{2pt}]
  \item A hard $3\times3$ reduces to three easy $2\times2$s you already know.
  \item Determinant $=$ signed volume: it turns a \emph{box} into a \emph{number}.
  \item Same story as 5.0, one dimension up: area $\to$ volume.
  \item $\det=0$ is the fingerprint of dependence: the box goes flat, the inverse dies.
\end{itemize}
\end{tabularx}
\end{skillbox}

\begin{skillbox}[{Vocabulary, Concepts, \& Theorems}]{greenbox}
\renewcommand{\arraystretch}{1.3}
\begin{tabularx}{\linewidth}{@{}l X@{}}
\textbf{Minor} & The $2\times2$ determinant left after crossing out an entry's row and column. \\
\textbf{Cofactor expansion} & $\det A$ across a row: each entry $\times$ its minor, with signs $+\,-\,+$. \\
\textbf{Signed volume} & $|\det A|$ is the volume of the box (parallelepiped) the three columns span; the \emph{sign} is orientation. \\
\textbf{Triangular shortcut} & If all entries below (or above) the diagonal are $0$, $\det A$ is the \emph{product of the diagonal}. \\
\textbf{Singular / invertible} (recall §2) & $A$ is invertible exactly when $\det A\ne0$; coplanar columns $\Rightarrow\det=0\Rightarrow$ no inverse. \\
\end{tabularx}
\end{skillbox}

\begin{fixedskillbox}[Activate Prior Knowledge \& Spiral Review]{blush}
    The Warm-Up rehearses the three pieces of cofactor expansion:
    \textbf{(1)} a \emph{$2\times2$ determinant} $ad-bc$ (Lesson~5.0) --- the minor students will
    compute three of; \textbf{(2)} the \emph{volume} of a rectangular box (length $\times$ width
    $\times$ height) --- the geometric meaning of the answer; and \textbf{(3)} the \emph{crossing-out
    motion} --- delete a row and column of a $3\times3$ and read off the $2\times2$ that stays. Debrief
    item~3 aloud: its answer is exactly the \emph{first} term of the notes' worked example.
\end{fixedskillbox}

\begin{skillbox}[Hook]{blush}
    In Lesson~5.0 two columns spanned a parallelogram and $|ad-bc|$ was its \emph{area}. Now put a
    \emph{third} column into space: the three columns span a slanted \textbf{box}. Ask:
    \textbf{``What one number measures how much space that box fills --- and how would you even compute
    it for a $3\times3$?''} Reveal the two answers together: geometrically it is the \emph{volume}
    ($|\det A|$); mechanically you don't need anything new --- a $3\times3$ determinant is just
    \emph{three $2\times2$ determinants you already know}, combined with the signs $+\,-\,+$. Name the
    payoff: the same ``one number reads off the geometry'' idea as 5.0, lifted one dimension. Pose the
    thread: \textbf{``When is that volume zero --- and what does a flat box say about the matrix?''}
\end{skillbox}

\begin{skillbox}[Lesson]{blush}
\begin{multicols}{2}
\textbf{1. From area to volume.} Two columns $\to$ parallelogram (area); \emph{three} columns $\to$ a
box in space (a parallelepiped). $|\det A|$ is its \textbf{volume} --- the 5.0 idea, one dimension up.
Show the box figure (three column-edges from a corner).

\textbf{2. Cofactor expansion.} Go across the top row. For each entry, cross out its row and column
(the \textbf{minor}, a $2\times2$) and attach signs $+\,-\,+$. Worked:
$A=\begin{bsmallmatrix} 2 & 1 & 0\\ 1 & 3 & 1\\ 0 & 1 & 2 \end{bsmallmatrix}$,
$\det A=2\!\cdot\!5-1\!\cdot\!2+0=8$. Stress: each piece is a Lesson~5.0 determinant, and the top-row
$0$ kills a term (expand where the zeros are).

\columnbreak

\textbf{3. Volume \& the triangular shortcut.} Axis box
$\mathrm{diag}(2,3,4)\to\det=24=2\!\cdot\!3\!\cdot\!4$. For a triangular matrix (like Unit~2's $U$),
$\det$ is the \textbf{product of the diagonal}: $\det\begin{bsmallmatrix} 2&5&1\\0&3&4\\0&0&2 \end{bsmallmatrix}=12$.

\textbf{4. When $\det=0$: collapse \& invertibility.} If the three columns lie in \emph{one plane}
(one is a combination of the others), the box is \emph{flat} --- zero volume --- so $\det=0$. E.g.\
$Z=\begin{bsmallmatrix} 1&0&1\\0&1&1\\1&1&2 \end{bsmallmatrix}$ (col$_3=$col$_1+$col$_2$): $\det=0$.
That is exactly when $A$ is \textbf{not invertible} (§2). \textbf{Road ahead:} \textbf{5.2} properties
of determinants; \textbf{5.3} transformations, where $|\det A|$ scales every volume.
\end{multicols}
\end{skillbox}

\begin{skillbox}[Active Monitoring]{redbox}
Circulate during the notes practice and Tier work. Watch for and correct:
\begin{itemize}[leftmargin=1.2em, nosep]
  \item the \emph{middle sign} --- the second term is $\mathbf{-}\,a_2\times(\text{its minor})$, not $+$;
  \item wrong \emph{minors} --- crossing out the wrong row/column, or an arithmetic slip inside the $2\times2$;
  \item forgetting \emph{absolute value} for volume, and reading $\det=0$ as ``no answer'' instead of ``coplanar / flat / no inverse.''
\end{itemize}
Cold-call prompts: ``Which entry did you cross out to get that $2\times2$?'' ``Why is the middle term
negative?'' ``A determinant of $0$ --- what do the three columns look like in space?''
\end{skillbox}

\begin{skillbox}[Group Work \& Differentiation]{redbox}
\textbf{Building Volume in 3D} activity (tiered; mirrors the notes).
\begin{multicols}{3}
\textbf{Tier R --- Remediate.} Expand a $3\times3$ (exploit a zero); find the volume of an axis box;
see that a zero row flattens the box ($\det=0$).

\columnbreak
\textbf{Tier A --- Approaching.} Confirm the triangular shortcut two ways; solve for the entry that
makes $\det=0$ and connect it to \emph{invertibility}; use the \emph{swap rule} for orientation.

\columnbreak
\textbf{Tier E --- Extension.} \emph{First look at 3D transformations:} apply a matrix to the unit
cube, find the image volume $=|\det A|$ (including a reflection with $\det<0$ and a $\det=0$ collapse).
\end{multicols}
\end{skillbox}

\begin{skillbox}[Individual Work \& Assessment]{redbox}
\textbf{Exit Ticket} (independent, no notes): compute a $3\times3$ determinant by expansion; find the
\emph{volume} of a box (choosing a zero-rich row); and a \textbf{justification check} --- three
coplanar columns give $\det=0$ and no inverse, explained in one sentence. Collect and sort into ``got
it / expansion slip / concept slip'' to decide whether 5.2 opens with a quick expansion re-warm.
\end{skillbox}

\begin{skillbox}[Reinforcement \& Extension]{goldbox}
\textbf{Homework:} compute $3\times3$ determinants (including a triangular one and a $\det=0$); volume
from three columns; solve-for-the-entry to force $\det=0$ and tie it to invertibility (§2); and a short
justification item. \textbf{Extension:} apply a matrix to the unit cube and show the image volume is
$|\det A|$ --- a first taste of Lesson~5.3, including a reflection. \textbf{Preview:} Lesson~5.2,
\emph{Properties and Applications of Determinants} --- the rules that let you find a determinant
\emph{without} always expanding.
\end{skillbox}
\end{document}
